\documentclass[11pt]{article}
\usepackage[margin=1in]{geometry}
\usepackage{amsmath,amssymb}
\usepackage[hidelinks]{hyperref}
\emergencystretch=3em

\title{Partial Literature Answer for a Basis-Reliance Question in arXiv:1504.06862}
\author{Run fa\_banach\_001, lane 7}
\date{June 14, 2026}

\begin{document}
\maketitle

\section*{Status}

This note records a partial later-literature answer to the basis-reliance
remark after Theorem 1.2 of Kurka, arXiv:1504.06862. The later paper is
Kurka, arXiv:1507.03899.

The answer is partial: it removes the input monotone-basis hypothesis for
analytic families of spaces with separable dual and for analytic families of
separable reflexive spaces. It does not settle every class appearing in the
source theorem, and it does not address the separate Schur-space issue from
Remark 2.4(c).

\section*{Source Signal}

The source paper proves an isometric amalgamation theorem for analytic classes
whose members have monotone bases. Immediately after the theorem, it says that
the author does not know whether the reliance on a basis can be dropped, as in
the isomorphic setting.

The theorem includes, among other classes, spaces with a monotone shrinking
basis and reflexive spaces with a monotone basis.

\section*{Later Partial Answer}

The later paper arXiv:1507.03899 explicitly returns to this obstruction. Its
introduction says that the technique from arXiv:1504.06862 applies only to
spaces with a monotone basis, so the basis-reliance problem appears again, and
that for two considered classes the problem can be solved by the same method as
in the isomorphic setting.

Its main theorem states:
\begin{enumerate}
\item If $\mathcal C$ is an analytic set of Banach spaces with separable dual,
then there is a Banach space $E$ with a shrinking monotone basis which contains
an isometric copy of every member of $\mathcal C$.
\item If $\mathcal C$ is an analytic set of separable reflexive Banach spaces,
then there is a reflexive Banach space $E$ with a monotone basis which contains
an isometric copy of every member of $\mathcal C$.
\end{enumerate}

Consequently, in these two branches the input-class basis hypothesis from
arXiv:1504.06862 is no longer needed. The later theorem gives isometric copies;
this note does not claim that the copies are $1$-complemented.

\section*{Scope}

This packet is a literature-status record, not a new theorem. It should be
indexed as a partial subcase answer. It does not resolve the source theorem's
non-isometrically-universal class without bases, the strictly convex class
without bases, or the Schur-tree question from Remark 2.4(c).

\section*{References}

\begin{enumerate}
\item O. Kurka, \emph{Amalgamations of classes of Banach spaces with a
monotone basis}, arXiv:1504.06862.
\item O. Kurka, \emph{Zippin's embedding theorem and amalgamations of classes
of Banach spaces}, arXiv:1507.03899.
\end{enumerate}

\end{document}
